Table~\ref{tab:skill} and Fig.~\ref{fig:lead} give the forecast skill of every system in the \texttt{main} group; Table~\ref{tab:clim} gives the free-run statistics. Values are means $\pm$ standard deviation over 5 seeds; differences and $p$-values come from \texttt{rh compare} (paired $t$-test over seeds, reference CNN $K{=}1$, so a positive $\Delta$ means lower RMSE than $K{=}1$).

\begin{table*}[t]\centering
\caption{Forecast skill, group \texttt{main}. RMSE at 1 step and at leads of 1, 2, 3 and 5 time units (t10, t20, t30, t50) and ACC lead time (time units until mean ACC $<0.6$); lower RMSE and higher lead time are better. Bold/underline mark best/second best per column. Climatology's ACC is zero by definition, so its ``lead time'' is an interpolation artefact.}\label{tab:skill}
\begin{tabular}{lccccccc}
\toprule
Method & rmse\_1step $\downarrow$ & rmse\_t10 $\downarrow$ & rmse\_t20 $\downarrow$ & rmse\_t30 $\downarrow$ & rmse\_t50 $\downarrow$ & acc\_leadtime $\uparrow$ & $n$ \\
\midrule
Persistence & 0.923 $\pm$ 0.005 & 5.160 $\pm$ 0.020 & 5.364 $\pm$ 0.014 & 4.948 $\pm$ 0.027 & 5.088 $\pm$ 0.031 & 0.190 $\pm$ 0.001 & 5 \\
Climatology & 3.538 $\pm$ 0.007 & 3.534 $\pm$ 0.008 & 3.540 $\pm$ 0.010 & 3.538 $\pm$ 0.006 & \textbf{3.536 $\pm$ 0.009} & 0.020 $\pm$ 0.000 & 5 \\
Linear stencil (ridge) & 0.698 $\pm$ 0.002 & 3.435 $\pm$ 0.036 & 3.732 $\pm$ 0.028 & 3.794 $\pm$ 0.021 & 3.753 $\pm$ 0.011 & 0.366 $\pm$ 0.009 & 5 \\
CNN K=1 & \textbf{0.023 $\pm$ 0.000} & 0.388 $\pm$ 0.015 & 1.308 $\pm$ 0.034 & 2.385 $\pm$ 0.038 & 3.775 $\pm$ 0.041 & 3.984 $\pm$ 0.038 & 5 \\
\textbf{CNN K=4} & \underline{0.024 $\pm$ 0.000} & \underline{0.374 $\pm$ 0.021} & \underline{1.290 $\pm$ 0.044} & \underline{2.359 $\pm$ 0.044} & 3.752 $\pm$ 0.035 & \underline{4.030 $\pm$ 0.050} & 5 \\
CNN K=8 & 0.025 $\pm$ 0.000 & \textbf{0.357 $\pm$ 0.024} & \textbf{1.243 $\pm$ 0.048} & \textbf{2.317 $\pm$ 0.037} & \underline{3.733 $\pm$ 0.037} & \textbf{4.053 $\pm$ 0.043} & 5 \\
\bottomrule
\end{tabular}

\end{table*}

\begin{table*}[t]\centering
\caption{ACC at leads 0.5, 1, 2, 3, 5 time units (acc\_t05\ldots acc\_t50), group \texttt{main}. The climatology forecast has zero anomaly, so its ACC is set to $0$.}\label{tab:acc}
\begin{tabular}{lcccccc}
\toprule
Method & acc\_t05 $\uparrow$ & acc\_t10 $\uparrow$ & acc\_t20 $\uparrow$ & acc\_t30 $\uparrow$ & acc\_t50 $\uparrow$ & $n$ \\
\midrule
Persistence & -0.176 $\pm$ 0.006 & -0.067 $\pm$ 0.007 & -0.153 $\pm$ 0.008 & 0.023 $\pm$ 0.010 & -0.036 $\pm$ 0.010 & 5 \\
Climatology & 0.000 $\pm$ 0.000 & 0.000 $\pm$ 0.000 & 0.000 $\pm$ 0.000 & 0.000 $\pm$ 0.000 & 0.000 $\pm$ 0.000 & 5 \\
Linear stencil (ridge) & 0.550 $\pm$ 0.006 & 0.408 $\pm$ 0.012 & 0.210 $\pm$ 0.011 & 0.112 $\pm$ 0.011 & 0.039 $\pm$ 0.010 & 5 \\
CNN K=1 & 0.999 $\pm$ 0.000 & 0.994 $\pm$ 0.000 & 0.934 $\pm$ 0.003 & 0.778 $\pm$ 0.006 & 0.436 $\pm$ 0.012 & 5 \\
\textbf{CNN K=4} & \underline{0.999 $\pm$ 0.000} & \underline{0.995 $\pm$ 0.001} & \underline{0.936 $\pm$ 0.004} & \underline{0.783 $\pm$ 0.007} & \underline{0.443 $\pm$ 0.011} & 5 \\
CNN K=8 & \textbf{0.999 $\pm$ 0.000} & \textbf{0.995 $\pm$ 0.001} & \textbf{0.940 $\pm$ 0.004} & \textbf{0.791 $\pm$ 0.006} & \textbf{0.449 $\pm$ 0.009} & 5 \\
\bottomrule
\end{tabular}

\end{table*}

\begin{figure}[t]\centering\includegraphics[width=\columnwidth]{lead_curves.pdf}
\caption{RMSE (log scale, left) and ACC (right) against lead time; lines are seed means, bands the min--max over seeds. The three CNN curves nearly coincide. The horizontal line is the ACC${=}0.6$ threshold.}\label{fig:lead}\end{figure}

\begin{table*}[t]\centering
\caption{Free-run statistics (200 time units from 16 states), group \texttt{main}. \texttt{var\_ratio}, \texttt{roughness\_ratio}: emulator/truth, target $1$; the bold/underline ranking in this table is by ``lower is better'' and should be ignored for these two columns (the linear stencil, whose runs decay to the mean, is ``best'' by that rule). Persistence and climatology have no free run.}\label{tab:clim}
\begin{tabular}{lcccccc}
\toprule
Method & stable\_frac $\uparrow$ & survival\_time $\uparrow$ & var\_ratio $\downarrow$ & roughness\_ratio $\downarrow$ & ks $\downarrow$ & $n$ \\
\midrule
Linear stencil (ridge) & \textbf{1.000 $\pm$ 0.000} & \textbf{200.000 $\pm$ 0.000} & \textbf{0.001 $\pm$ 0.000} & \textbf{0.001 $\pm$ 0.000} & 0.468 $\pm$ 0.002 & 5 \\
CNN K=1 & \underline{1.000 $\pm$ 0.000} & \underline{200.000 $\pm$ 0.000} & 1.003 $\pm$ 0.003 & 1.005 $\pm$ 0.007 & 0.002 $\pm$ 0.001 & 5 \\
\textbf{CNN K=4} & 1.000 $\pm$ 0.000 & 200.000 $\pm$ 0.000 & 1.002 $\pm$ 0.004 & 1.005 $\pm$ 0.010 & \underline{0.002 $\pm$ 0.002} & 5 \\
CNN K=8 & 1.000 $\pm$ 0.000 & 200.000 $\pm$ 0.000 & \underline{0.998 $\pm$ 0.003} & \underline{0.997 $\pm$ 0.006} & \textbf{0.002 $\pm$ 0.001} & 5 \\
\bottomrule
\end{tabular}

\end{table*}

\subsection{H4: the emulators have skill far beyond the baselines}
Persistence has an ACC lead time of only \rhval{main/persistence/l96/acc_leadtime/mean:2} time units, because the L96 waves propagate and persistence ACC oscillates around zero (Fig.~\ref{fig:lead}); the linear stencil reaches \rhval{main/linear-stencil-ridge/l96/acc_leadtime/mean:2}, and the CNNs between \rhval{main/cnn-k-1/l96/acc_leadtime/mean:2} ($K{=}1$) and \rhval{main/cnn-k-8/l96/acc_leadtime/mean:2} ($K{=}8$). The one-step RMSE of $K{=}1$ is \rhval{main/cnn-k-1/l96/rmse_1step/mean:3} against \rhval{main/linear-stencil-ridge/l96/rmse_1step/mean:3} for the linear stencil, which is hardly better than persistence's \rhval{main/persistence/l96/rmse_1step/mean:3}, so the nonlinearity matters. H4 is \textbf{supported}: e.g.\ the lead-time difference between $K{=}1$ and persistence is \rhval{cmp/main/persistence/l96/acc_leadtime/delta:2} time units (paired $p=\rhval{cmp/main/persistence/l96/acc_leadtime/paired_p}$).

\textit{Where it fails.} At 5 time units the CNN RMSE (\rhval{main/cnn-k-1/l96/rmse_t50/mean:2} for $K{=}1$) is \emph{above} climatology's \rhval{main/climatology/l96/rmse_t50/mean:2}, although its ACC is still \rhval{main/cnn-k-1/l96/acc_t50/mean:2}: the emulator is sharp (it does not blur toward the mean, see H3), and a plausible reading is that once phase errors grow a sharp forecast is penalised by RMSE more than the constant mean is (we did not test this with a blurred model). Judged by RMSE alone, the CNN beats climatology only before this crossing; at 3 time units it is still better (\rhval{main/cnn-k-1/l96/rmse_t30/mean:2} against \rhval{main/climatology/l96/rmse_t30/mean:2}). At 5 time units the linear stencil (\rhval{main/linear-stencil-ridge/l96/rmse_t50/mean:2}) is statistically indistinguishable from $K{=}1$ in RMSE (paired $p=\rhval{cmp/main/linear-stencil-ridge/l96/rmse_t50/paired_p:2}$), and both are worse than climatology; the linear model reaches this level by decaying to the mean, not by being skilful.

\subsection{H1: rollout training helps modestly at medium range}
RMSE at lead 2 falls monotonically in the mean with $K$: \rhval{main/cnn-k-1/l96/rmse_t20/mean:3} ($K{=}1$), \rhval{sweep_k/cnn-k-2/l96/rmse_t20/mean:3} ($K{=}2$), \rhval{main/cnn-k-4/l96/rmse_t20/mean:3} ($K{=}4$), \rhval{main/cnn-k-8/l96/rmse_t20/mean:3} ($K{=}8$) (Fig.~\ref{fig:sweepk}). Against $K{=}1$, $K{=}8$ improves by \rhval{cmp/main/cnn-k-8/l96/rmse_t20/delta:3} ($K{=}1$ is \rhval{cmp/main/cnn-k-8/l96/rmse_t20/rel_delta:pct1}\% higher than $K{=}8$; paired $p=\rhval{cmp/main/cnn-k-8/l96/rmse_t20/paired_p:3}$, Cohen's $d=\rhval{cmp/main/cnn-k-8/l96/rmse_t20/cohen_d:2}$; Welch $p=\rhval{cmp/main/cnn-k-8/l96/rmse_t20/welch_p:2}$), and $K{=}4$ by \rhval{cmp/main/cnn-k-4/l96/rmse_t20/delta:3} ($K{=}1$ is \rhval{cmp/main/cnn-k-4/l96/rmse_t20/rel_delta:pct1}\% higher than $K{=}4$; paired $p=\rhval{cmp/main/cnn-k-4/l96/rmse_t20/paired_p:3}$, but Welch $p=\rhval{cmp/main/cnn-k-4/l96/rmse_t20/welch_p:2}$). The same ordering holds at leads 1, 3 and 5 (Table~\ref{tab:skill}) and for the ACC lead time, which rises from \rhval{main/cnn-k-1/l96/acc_leadtime/mean:3} ($K{=}1$) to \rhval{main/cnn-k-8/l96/acc_leadtime/mean:3} ($K{=}8$; paired $p=\rhval{cmp/main/cnn-k-8/l96/acc_leadtime/paired_p:3}$). H1 is therefore \textbf{supported for $K{=}8$ and only weakly for $K{=}4$}: the effect is small (about a twentieth of the error at $K{=}8$, below the relative gain we had set as a meaningful threshold in the study configuration), the $K{=}4$ evidence rests on the paired test alone, no multiplicity correction is applied, and seed-to-seed variation (std \rhval{main/cnn-k-1/l96/rmse_t20/std:3}) is a sizeable fraction of the effect. At the longest lead 5, $K{=}4$ is not distinguishable from $K{=}1$ (paired $p=\rhval{cmp/main/cnn-k-4/l96/rmse_t50/paired_p:2}$), whereas $K{=}8$ is (paired $p=\rhval{cmp/main/cnn-k-8/l96/rmse_t50/paired_p:3}$), though the absolute gain is only \rhval{cmp/main/cnn-k-8/l96/rmse_t50/delta:2}.

\textit{Cost.} The one-step RMSE gets worse with $K$ (\rhval{main/cnn-k-1/l96/rmse_1step/mean:4}, \rhval{sweep_k/cnn-k-2/l96/rmse_1step/mean:4}, \rhval{main/cnn-k-4/l96/rmse_1step/mean:4}, \rhval{main/cnn-k-8/l96/rmse_1step/mean:4} for $K=1,2,4,8$; $K{=}8$ vs $K{=}1$ paired $p=\rhval{cmp/main/cnn-k-8/l96/rmse_1step/paired_p}$), and the training compute per step grows about linearly in $K$.

\begin{figure}[t]\centering\includegraphics[width=\columnwidth]{sweep_K_panels.pdf}
\caption{Effect of the training rollout length $K$ (groups \texttt{main} and \texttt{sweep\_K}): one-step RMSE, RMSE at lead 2 and ACC lead time. Grey lines are the five seeds (data seeds are shared across $K$), red is the seed mean.}\label{fig:sweepk}\end{figure}

\subsection{H2: no instability was observed}
All free runs of 200 time units survive for the linear stencil and for every CNN: \texttt{stable\_frac} is \rhval{main/cnn-k-1/l96/stable_frac/mean:1} for $K=1$ and for every other $K$ (Table~\ref{tab:clim}). H2 is \textbf{refuted in this setting}: there is nothing for multi-step training to fix, and \texttt{rh compare} has no variance to test. We only observed $K{=}1$ with $\Delta{=}0.05$; other steps, architectures or noisier data might behave differently.

\subsection{H3: climatological variance is preserved, not blurred}
The variance ratio of the CNNs is \rhval{main/cnn-k-1/l96/var_ratio/mean:3} ($K{=}1$), \rhval{main/cnn-k-4/l96/var_ratio/mean:3} ($K{=}4$) and \rhval{main/cnn-k-8/l96/var_ratio/mean:3} ($K{=}8$); the roughness ratio is \rhval{main/cnn-k-1/l96/roughness_ratio/mean:3}, \rhval{main/cnn-k-4/l96/roughness_ratio/mean:3} and \rhval{main/cnn-k-8/l96/roughness_ratio/mean:3}, and the KS distance is at most \rhval{main/cnn-k-1/l96/ks/mean:3} for $K{=}1$. The difference $K{=}8$ vs $K{=}1$ is \rhval{cmp/main/cnn-k-8/l96/var_ratio/delta:3} (paired $p=\rhval{cmp/main/cnn-k-8/l96/var_ratio/paired_p:2}$, Welch $p=\rhval{cmp/main/cnn-k-8/l96/var_ratio/welch_p:3}$) and $K{=}4$ vs $K{=}1$ is \rhval{cmp/main/cnn-k-4/l96/var_ratio/delta:3} ($p=\rhval{cmp/main/cnn-k-4/l96/var_ratio/paired_p:2}$). H3 is \textbf{not supported}: the sign is the predicted one for $K{=}8$ but the effect is below one percent and not significant under the registered paired test (the Welch test alone gives a smaller $p$, but the effect is below one percent and we do not rest a claim on it), so we cannot claim multi-step training reduces variance. By contrast the ridge linear stencil is ``stable'' only in the trivial sense: its free runs decay to the mean (variance ratio \rhval{main/linear-stencil-ridge/l96/var_ratio/mean:4}, KS \rhval{main/linear-stencil-ridge/l96/ks/mean:2}), which shows that the stability criterion must be read together with the variance ratio.
